% ID: MAT-TRI-EQG-009
% Titolo: Equazioni goniometriche con fattorizzazione (2)
% Disciplina: Matematica
% Area: Trigonometria
% Argomento: Equazioni goniometriche
% Classe: Quarta secondaria di secondo grado
% Difficolta: 4
% Tipo: problema_numerico
% Tag: matematica, trigonometria, equazioni_goniometriche
% Fonte: verifica_goniometria_intermedia_anonimizzata
% Provenienza: materiale fornito dal docente; testo anonimizzato; soluzione redatta per Exergo

\begin{esercizio}
Risolvi le seguenti equazioni: 

\begin{enumerate}
\item $\displaystyle \sin x+\sqrt{3}\cos x=1;$

\item $\displaystyle \dfrac{\sin 2x}{1+\sin x}=\cos x.$
\end{enumerate}
Salvo diversa indicazione, lavora in $\mathbb R$; angoli in radianti.
\end{esercizio}

\begin{soluzione}
a) $2\sin(x+\pi/3)=1$ dà $x=-\pi/6+2k\pi$ o $x=\pi/2+2k\pi$.
b) Il dominio esclude $\sin x=-1$. Si ottiene $\cos x(\sin x-1)=0$.
Resta soltanto $x=\pi/2+2k\pi$. In tutte le famiglie di soluzioni $k\in\mathbb{Z}$.
\end{soluzione}
